\chapter{Agents IA et propositions}
\label{ch:agents}

mybibli peut s'ouvrir à un agent d'intelligence artificielle ---
Claude Code, Hermes Agent ou tout autre client du \emph{Model Context
Protocol}\index{MCP}\index{Model Context Protocol} (MCP). L'agent lit
votre catalogue et suggère des améliorations : une série et ses tomes
manquants, une description ou un éditeur absents, une couverture.
Chaque suggestion arrive sous la forme d'une
\emph{proposition}\index{proposition} qui attend sur la page
\textbf{Propositions} qu'un bibliothécaire l'accepte, la modifie ou
la refuse.

Ce chapitre couvre :

\begin{itemize}
  \item ce qu'un agent peut faire, et ce qu'il ne peut pas faire,
  \item la création d'une clé \emph{propositions} et l'autorisation
        de l'adresse de l'agent,
  \item la connexion d'un client MCP --- d'abord les éléments
        génériques, puis Claude Code,
  \item les outils que l'agent voit,
  \item l'examen des propositions : liste, détail, comparaison,
        modification, acceptation, refus,
  \item les tomes attendus et la wish list,
  \item où attendent les couvertures proposées, et comment elles
        sont nettoyées.
\end{itemize}

\section{Ce qu'un agent peut faire, et ce qu'il ne peut pas faire}
\label{sec:agents-scope}

\begin{note}
\textbf{Un agent n'écrit jamais dans le catalogue.} Il lit les
titres, les séries et les contributeurs, et il dépose des
propositions. Rien ne change dans le catalogue tant qu'un
Bibliothécaire ou un Admin n'accepte pas une proposition sur la page
Propositions. Refuser une proposition, ou simplement l'ignorer, laisse
le catalogue exactement tel qu'il était.
\end{note}

Avec une clé \emph{propositions}, un agent peut :

\begin{itemize}
  \item chercher et lire les titres (champs bibliographiques,
        contributeurs, volumes et leur emplacement, positions dans
        les séries), les contributeurs et les séries --- y compris
        les lacunes d'une série fermée ;
  \item lister les titres sans couverture, ou aux métadonnées
        incomplètes ;
  \item proposer une nouvelle série, ou des ajouts à une série
        existante (des titres que vous possédez, et des \emph{tomes
        attendus} que vous n'avez pas encore) ;
  \item proposer une correction du titre principal (orthographe,
        ponctuation, ou le vrai titre d'une fiche qui ne porte que son
        ISBN), et des valeurs pour le sous-titre, la description,
        l'éditeur, l'année de publication, le code Dewey d'un titre,
        et des contributeurs à ajouter ;
  \item proposer une image de couverture ;
  \item consulter l'historique de ses propres propositions, avec vos
        décisions et les motifs de vos refus.
\end{itemize}

Il ne peut pas :

\begin{itemize}
  \item créer, modifier ou supprimer quoi que ce soit dans le
        catalogue --- titres, volumes, séries, contributeurs,
        emplacements, prêts ;
  \item voir les emprunteurs, les prêts, les utilisateurs, les
        réglages ni les propositions d'une autre clé ;
  \item effacer une valeur : une proposition complète ou corrige un
        champ, elle ne le vide jamais, et elle ne fait
        qu'\emph{ajouter} des contributeurs ;
  \item faire télécharger quoi que ce soit au serveur : une
        couverture proposée est envoyée sous forme d'image, et les
        adresses web que l'agent cite comme sources vous sont
        présentées comme liens, jamais consultées par mybibli.
\end{itemize}

Chaque proposition doit citer ses \textbf{sources} --- les pages
\texttt{http} ou \texttt{https} où l'agent a trouvé l'information ---
pour que vous puissiez les vérifier avant de décider.

\section{Avant de commencer : le réseau}
\label{sec:agents-network}

Le point d'entrée MCP se trouve à
\texttt{/api/mcp}\index{api/mcp@\texttt{/api/mcp}}, sur le même
serveur web que le reste de mybibli. Il est prévu pour un agent qui
tourne sur un ordinateur \emph{de votre réseau local}. Comme le reste
de mybibli (chapitre~\ref{ch:installation}), il ne doit jamais être
publié sur Internet : la clé circule dans chaque requête, en clair
sur un réseau local en HTTP simple.

\section{Étape 1 --- créer une clé \emph{propositions}}
\label{sec:agents-key}

\begin{enumerate}
  \item En tant qu'Admin, ouvrez \textbf{Administration
        \(\rightarrow\) Clés d'API} (section~\ref{sec:api-keys-create}).
  \item Donnez à la clé un libellé qui nomme l'agent, par ex.\
        \emph{« Claude Code sur le portable du bureau »}. Le libellé
        s'affiche à côté de chaque proposition que la clé dépose.
  \item Choisissez la portée \textbf{Lecture +
        propositions}\index{clé propositions}.
  \item Cliquez sur \textbf{Créer la clé} et copiez-la depuis la
        boîte de dialogue : elle n'est affichée qu'une fois. Une clé
        propositions commence par \texttt{mybibli\_pr\_}.
\end{enumerate}

Seule une clé propositions ouvre \texttt{/api/mcp} ; une clé en
lecture seule ou en lecture + écriture y est refusée. Utilisez une clé
par agent : la révoquer (section~\ref{sec:agents-revoke}) arrête alors
exactement cet agent.

\section{Étape 2 --- autoriser le nom sous lequel l'agent appelle mybibli}
\label{sec:agents-hosts}

Pour se défendre du \emph{DNS rebinding} (une page web qui amène un
navigateur à parler à votre NAS sous un nom étranger), le point
d'entrée ne répond qu'aux requêtes dont l'en-tête \texttt{Host}
figure sur une liste. Cet en-tête porte le nom (ou l'adresse) du
\emph{serveur} tel que l'agent l'a écrit dans son adresse --- pour
\texttt{http://nas.local:8080/api/mcp}, c'est
\texttt{nas.local:8080}. Il ne dit rien de la machine sur laquelle
tourne l'agent : la liste ne filtre pas les clients, tous les
ordinateurs de votre réseau qui appellent la même adresse envoient le
même \texttt{Host}, et un sous-réseau n'y aurait aucun sens. C'est la
clé propositions qui contrôle l'accès. Deux variables d'environnement, lues au
démarrage, la règlent (chapitre~\ref{ch:installation}, tableau de la
section~\ref{sec:env-vars}) :

\begin{description}
  \item[\texttt{MYBIBLI\_MCP\_ALLOWED\_HOSTS}]\index{variables d'env!MYBIBLI\_MCP\_ALLOWED\_*}
        Noms d'hôte ou adresses séparés par des virgules, chacun avec
        ou sans port (\texttt{hôte} ou \texttt{hôte:port} ; une
        entrée sans port vaut pour tous les ports). Par défaut, non
        définie ou vide : \texttt{localhost,127.0.0.1,::1} --- ce qui
        ne suffit que si l'agent tourne sur la machine qui héberge
        mybibli. Une valeur définie \emph{remplace} celle par
        défaut : sur un NAS, listez chaque nom et chaque adresse sous
        lesquels les agents atteignent le NAS (son nom d'hôte, son
        adresse IP si un agent l'utilise), et gardez les noms locaux
        si un client local en a besoin :
\begin{lstlisting}
MYBIBLI_MCP_ALLOWED_HOSTS=nas.local,192.168.1.10,localhost
\end{lstlisting}
  \item[\texttt{MYBIBLI\_MCP\_ALLOWED\_ORIGINS}]\index{variables d'env!MYBIBLI\_MCP\_ALLOWED\_*}
        Origines séparées par des virgules (\texttt{schéma://hôte:port},
        \texttt{:*} pour tout port). Par défaut : vide. Les clients
        MCP comme Claude Code n'envoient pas d'en-tête
        \texttt{Origin} et n'ont besoin de rien ici ; une requête qui
        porte un \texttt{Origin} est refusée s'il ne figure pas dans
        la liste. N'en listez une que pour un client ou un proxy qui
        en transmet une (le MCP Inspector,
        section~\ref{sec:agents-inspector}). Les pages web ne peuvent
        pas du tout appeler le point d'entrée : il n'envoie aucun
        en-tête CORS.
\end{description}

Définissez-les dans \texttt{.env} (le \texttt{docker-compose.yml}
fourni transmet les deux) et redémarrez le conteneur. Ce sont des
réglages de déploiement, et non des réglages d'administration, à
dessein : une page de l'interface web ne doit pas pouvoir élargir une
défense réseau.

\section{Étape 3 --- connecter l'agent}
\label{sec:agents-connect}

\subsection{Ce dont tout client MCP a besoin}
\label{sec:agents-connection-facts}

Quel que soit le client --- Claude Code, Hermes Agent ou un autre ---,
il lui faut ces éléments :

\begin{center}\small
\begin{tabular}{@{} l >{\raggedright\arraybackslash}p{9.5cm} @{}}
\toprule
\textbf{Réglage} & \textbf{Valeur} \\
\midrule
Adresse   & \texttt{http://<hôte>:<port>/api/mcp}, où \texttt{<hôte>}
            est l'un des hôtes autorisés de la
            section~\ref{sec:agents-hosts} et \texttt{<port>} le port
            sur lequel vous ouvrez mybibli (\texttt{HOST\_PORT}, 8080
            par défaut). \\
Transport & \emph{Streamable HTTP} (parfois simplement appelé
            \emph{HTTP}). Le serveur est sans état : chaque appel est
            un \texttt{POST} auquel il répond en JSON ; pas de flux
            d'événements, pas de session à conserver. \\
En-tête   & \texttt{Authorization: Bearer <clé propositions>} sur
            chaque requête. \\
Portée    & La clé doit être une clé \emph{propositions}
            (section~\ref{sec:agents-key}). Ni OAuth, ni connexion. \\
Protocole & Révisions MCP \texttt{2025-11-25} et
            \texttt{2026-07-28}. Les révisions plus anciennes sont
            refusées. \\
\bottomrule
\end{tabular}
\end{center}

Un client capable de poser un en-tête personnalisé sur un serveur
streamable HTTP peut se connecter. Les clients uniquement stdio et
l'ancien transport « HTTP + SSE » ne sont pas pris en charge.

\subsection{Claude Code}
\label{sec:agents-claude-code}

Enregistrez le serveur une fois pour toutes avec \texttt{claude mcp
add}. Remplacez l'adresse et la clé par les vôtres :

\begin{lstlisting}
claude mcp add --transport http mybibli \
  http://nas.local:8080/api/mcp \
  --header "Authorization: Bearer mybibli_pr_..."
\end{lstlisting}

La portée par défaut (\texttt{-s local}) garde le serveur --- et sa
clé --- dans votre propre configuration de Claude Code, pour le
dossier courant ; \texttt{-s user} le rend disponible dans tous les
dossiers. Évitez \texttt{-s project} : la clé serait écrite dans un
fichier \texttt{.mcp.json} destiné à être partagé avec un dépôt.

Vérifiez avec \texttt{claude mcp list} (le serveur doit apparaître
connecté), ou avec \texttt{/mcp} dans une session. Claude voit les
outils sous le nom \texttt{mcp\_\_mybibli\_\_<outil>}, par ex.\
\texttt{mcp\_\_mybibli\_\_search\_titles}.

\subsubsection*{Sans enregistrement : \texttt{--mcp-config}}

Pour une session ponctuelle ou un script, décrivez le serveur dans un
fichier JSON et passez-le sur la ligne de commande. Gardez ce fichier
privé (il contient la clé), par ex.\ \texttt{chmod 600
mybibli-mcp.json} :

\begin{lstlisting}
{
  "mcpServers": {
    "mybibli": {
      "type": "http",
      "url": "http://nas.local:8080/api/mcp",
      "headers": { "Authorization": "Bearer mybibli_pr_..." }
    }
  }
}
\end{lstlisting}

\begin{lstlisting}
claude --mcp-config mybibli-mcp.json
# non interactif, outils mybibli seulement :
claude -p "Quelles séries fermées ont des lacunes ?" \
  --strict-mcp-config --mcp-config mybibli-mcp.json \
  --allowedTools "mcp__mybibli__list_series"
\end{lstlisting}

\subsection{Vérifier le point d'entrée avec le MCP Inspector}
\label{sec:agents-inspector}

Le MCP Inspector est le client de test de référence du protocole.
C'est le moyen le plus rapide de distinguer un problème de réseau ou
de clé d'un problème d'agent. Il tourne avec Node.js (\texttt{npx} le
télécharge au premier usage) sur n'importe quelle machine qui atteint
mybibli. Dans les commandes ci-dessous, remplacez l'adresse et la
clé.

\begin{enumerate}
  \item \textbf{Lister les outils} (mode ligne de commande) :
\begin{lstlisting}
npx @modelcontextprotocol/inspector --cli \
  http://nas.local:8080/api/mcp --transport http \
  --header "Authorization: Bearer mybibli_pr_..." \
  --method tools/list
\end{lstlisting}
        Attendu : les douze outils de la
        section~\ref{sec:agents-tools}.
  \item \textbf{En appeler un} :
\begin{lstlisting}
npx @modelcontextprotocol/inspector --cli \
  http://nas.local:8080/api/mcp --transport http \
  --header "Authorization: Bearer mybibli_pr_..." \
  --method tools/call --tool-name get_title --tool-arg id=1
\end{lstlisting}
        Attendu : la fiche du titre~1 (ou une erreur si cet
        identifiant n'existe pas --- toute réponse de l'outil prouve
        la connexion).
  \item \textbf{Facultatif, l'interface web} : lancez
\begin{lstlisting}
npx @modelcontextprotocol/inspector
\end{lstlisting}
        sans argument, ouvrez l'adresse qu'il affiche, choisissez le
        transport \emph{Streamable HTTP}, saisissez l'adresse, ajoutez
        l'en-tête \texttt{Authorization}, puis \emph{Connect},
        \emph{List Tools}, et lancez \texttt{get\_title}. L'Inspector
        parle à mybibli au travers de son propre proxy local ; si ce
        proxy transmet son \texttt{Origin} et que la connexion est
        refusée avec un \emph{403}, ajoutez l'origine qu'il affiche
        (en général \texttt{http://localhost:6274}) à
        \texttt{MYBIBLI\_MCP\_ALLOWED\_ORIGINS}, redémarrez, puis
        retirez-la une fois la vérification faite.
\end{enumerate}

Chaque appel d'outil est journalisé au niveau \texttt{info} par une
ligne \texttt{MCP tool call} qui nomme l'outil et l'identifiant de la
clé (chapitre~\ref{ch:operations-debugging}). Les messages de la
bibliothèque MCP elle-même sont limités au niveau \texttt{warn}. Pour
les voir pendant une enquête, saisissez \texttt{info,rmcp=debug} comme
\textbf{Niveau de journalisation} dans \textbf{Administration
\(\rightarrow\) Système \(\rightarrow\) Journalisation} (effet
immédiat, sans redémarrage), puis remettez la valeur précédente.
Définir \texttt{MYBIBLI\_LOG\_LEVEL} dans \texttt{.env} ne sert à
rien ici : le niveau enregistré dans le panneau l'emporte à chaque
démarrage.

\subsection{Quand la connexion échoue}

\begin{center}\small
\begin{tabular}{@{} l >{\raggedright\arraybackslash}p{10cm} @{}}
\toprule
\textbf{Réponse} & \textbf{Cause probable} \\
\midrule
401 & Pas d'en-tête \texttt{Authorization}, une clé mal recopiée ou
      révoquée. \\
403 \texttt{insufficient\_scope} & La clé est en lecture seule ou en
      lecture + écriture ; créez une clé propositions. \\
403, autre & Le \texttt{Host} utilisé par le client ne figure pas
      dans \texttt{MYBIBLI\_MCP\_ALLOWED\_HOSTS}, ou la requête porte
      un \texttt{Origin} absent de
      \texttt{MYBIBLI\_MCP\_ALLOWED\_ORIGINS}. \\
400 / -32022 & (La seconde est une erreur JSON-RPC,
      \emph{Unsupported protocol version}.) Le client parle une révision du protocole que mybibli ne sert pas ; mettez
      le client à jour. \\
406 / 415 & Le client ne parle pas streamable HTTP (\texttt{Accept}
      ou \texttt{Content-Type} inadéquat). \\
405 & Le client a tenté un \texttt{GET} ou un \texttt{DELETE} (un
      ancien client HTTP + SSE) ; choisissez le transport streamable
      HTTP. \\
413 / 408 & Une requête de plus de $\approx$\,3,7\,Mio (couverture
      trop lourde) ou de plus de 30\,s. \\
\bottomrule
\end{tabular}
\end{center}

\section{Les outils}
\label{sec:agents-tools}

L'agent découvre les outils de lui-même ; chacun porte une
description rédigée à son intention. À la connexion, mybibli lui
indique aussi les sources qu'il interroge automatiquement : les
fournisseurs de métadonnées de chaque type de média dans leur ordre (la
recherche s'arrête au premier qui connaît le titre), et les sources de
couvertures essayées pour les livres et bandes dessinées dotés d'un
ISBN. L'agent est invité à chercher d'abord ailleurs, et prévenu que
ces sources enrichissent leurs catalogues et peuvent détenir une notice
aujourd'hui. Un fournisseur dont la clé
n'est pas configurée n'est pas cité. Pour mémoire :

\subsubsection*{Lecture}

\begin{description}[style=nextline]
  \item[\texttt{search\_titles}] Recherche en texte libre, ou par ISBN
        ou étiquette de volume (code V) exacts ; filtres par type de
        média et par genre.
  \item[\texttt{get\_title}] Un titre complet : champs, contributeurs
        et rôles, volumes et leur emplacement, positions dans les
        séries, présence d'une couverture, champs modifiés à la main.
  \item[\texttt{search\_contributors}] Les contributeurs par nom, avec
        le nombre de titres de chacun.
  \item[\texttt{list\_series}] Les séries par nom, avec les positions
        possédées et le nombre de lacunes ; peut se limiter aux
        séries lacunaires.
  \item[\texttt{get\_series}] Une série : ses positions, celles qui
        manquent à une série fermée, ses tomes attendus.
  \item[\texttt{list\_titles\_without\_cover}] Les titres sans image
        de couverture.
  \item[\texttt{list\_titles\_with\_incomplete\_metadata}] Les titres
        sans description, éditeur, contributeur ou année de
        publication.
  \item[\texttt{list\_my\_proposals}] Les propositions déposées avec
        \emph{cette} clé : statut, date de décision, motif de votre
        refus, et date à partir de laquelle une proposition
        équivalente peut être déposée à nouveau.
\end{description}

\subsubsection*{Propositions}

\begin{description}[style=nextline]
  \item[\texttt{propose\_series\_create}] Une nouvelle série --- nom,
        fermée (avec son nombre de volumes) ou ouverte --- avec les
        titres que vous possédez à leur position, et des tomes
        attendus.
  \item[\texttt{propose\_series\_attach}] Des titres à placer dans une
        série existante, et/ou des tomes attendus à lui ajouter.
  \item[\texttt{propose\_title\_metadata}] Une correction du titre, le sous-titre, la description,
        éditeur, année de publication, code Dewey et contributeurs à
        ajouter, pour un titre.
  \item[\texttt{propose\_cover}] Une image de couverture pour un titre
        (JPEG, PNG, WebP ou GIF, 2\,Mio au plus).
\end{description}

Un outil refuse, avec une explication sur laquelle l'agent peut agir,
une proposition qui ne pourrait jamais être acceptée : un titre
inconnu, une position déjà prise, une valeur qu'un formulaire de
bibliothécaire refuserait, une source manquante, une seconde
proposition en attente pour la même série, ou une proposition que
vous avez refusée récemment (section~\ref{sec:agents-delay}). Chaque
agent peut avoir au plus 100 propositions de couverture en attente à
la fois.

\section{Examiner les propositions}
\label{sec:proposals}

\subsection{Voir que des propositions attendent}

Les Bibliothécaires et les Admins voient une entrée
\textbf{Propositions}\index{page Propositions} dans la barre de
navigation. Tant que des propositions sont en attente, elle porte un
badge avec leur nombre (\texttt{99+} au-delà de 99), et la section
d'attention de la page d'accueil affiche un indicateur
\textbf{Propositions en attente}. Les deux mènent à
\texttt{/proposals}. Les visiteurs anonymes ne voient ni l'un ni
l'autre.

\subsection{La liste}

\texttt{/proposals} affiche par défaut les propositions en attente,
les plus récentes d'abord, 25 par page. Deux filtres la restreignent :

\begin{itemize}
  \item \textbf{Statut} --- \emph{En attente}, \emph{Acceptée},
        \emph{Refusée}, \emph{Périmée}, ou tous ;
  \item \textbf{Type} --- \emph{Nouvelle série}, \emph{Ajouts à une
        série}, \emph{Métadonnées}, \emph{Couverture}, ou tous.
\end{itemize}

Chaque ligne nomme l'objet --- le nom de série proposé, ou la série
ou le titre concerné --- et \textbf{l'objet ouvre la proposition}.
Un petit lien \textbf{Voir le titre} (ou \textbf{Voir la série}) à
côté mène à la fiche du catalogue. Sous l'objet, une proposition de
métadonnées liste les champs qu'elle propose (\emph{Proposé : Éditeur,
Année de publication}) ; à côté de l'objet, une proposition de
couverture en attente montre une miniature de l'image proposée. La ligne indique aussi la date de dépôt,
la clé, le statut et le nombre de lignes.

\subsection{La page de détail}

La page de détail indique qui a déposé la proposition et quand, son
statut, les \textbf{Sources} citées par l'agent (des liens qui
s'ouvrent dans un nouvel onglet --- vérifiez-les), et les
\textbf{Changements proposés}, une carte par ligne :

\begin{itemize}
  \item \textbf{Propositions de série} --- la ligne de série (nom,
        type, nombre de volumes), puis chaque titre avec sa position
        (et les positions que couvre une intégrale), puis chaque tome
        attendu avec sa position, son titre et son ISBN.
  \item \textbf{Propositions de métadonnées} --- un tableau de
        comparaison par champ, \textbf{Actuel} à côté de
        \textbf{Proposé} (pour les contributeurs : \textbf{À
        ajouter}), côte à côte même sur un téléphone. Un badge en
        toutes lettres indique si le champ est \emph{Nouveau} (vide
        aujourd'hui), \emph{Modifié} ou \emph{Inchangé} ; la valeur
        proposée d'un champ nouveau ou modifié est aussi marquée d'une
        bordure épaisse, si bien que la mise en évidence ne repose
        jamais sur la seule couleur.
  \item \textbf{Propositions de couverture} --- la couverture actuelle
        (ou l'image par défaut) à côté de l'image proposée, et une
        phrase qui dit si accepter \emph{ajoute} une couverture ou
        \emph{remplace} l'actuelle.
\end{itemize}

\subsection{Modifier avant d'accepter}

Tant qu'une proposition est en attente, vous pouvez la corriger
plutôt que la refuser. \textbf{Modifier} ouvre une ligne sur place ;
\textbf{Enregistrer} la vérifie avec les mêmes règles qu'au dépôt
par l'agent.

\begin{itemize}
  \item Une nouvelle série : son nom, son type et son nombre de
        volumes.
  \item Une ligne de titre : sa position, la case intégrale et la
        dernière position.
  \item Un tome attendu : sa position, son titre et son ISBN.
  \item Un champ de métadonnées : la valeur proposée (la valeur
        actuelle sert de référence et ne change pas) ; pour les
        contributeurs, chaque nom et rôle --- videz un nom pour
        écarter ce contributeur.
\end{itemize}

\textbf{Retirer} enlève une ligne de la proposition (titre, tome
attendu, champ de métadonnées ou ligne de contributeurs). La ligne
n'est pas appliquée à l'acceptation, mais elle reste consignée dans
une section \emph{Lignes retirées avant la décision}. Certaines lignes
ne peuvent pas être retirées : la ligne de série, l'unique ligne d'une
proposition de couverture et la dernière ligne d'une proposition de
métadonnées --- refusez alors la proposition. Une couverture ne se
modifie pas : une autre image est une nouvelle proposition.

Une ligne modifiée garde sous elle la valeur de l'agent
(\emph{L'assistant proposait : \dots}) et indique qui l'a modifiée et
quand. L'historique de l'agent montre toujours ce qu'il a déposé.

\subsection{Accepter}

\textbf{Accepter} ouvre une confirmation qui résume ce qui sera
écrit --- la série créée ou complétée et le nombre de titres et de
tomes attendus, les champs écrits sur le titre, les contributeurs
crédités, ou la couverture installée. Une fois confirmé :

\begin{itemize}
  \item une proposition de série crée la série, place les titres et
        ajoute les tomes attendus ; un tome attendu dont la série
        possède déjà le livre à une autre position est écarté, et le
        message le signale ;
  \item une proposition de métadonnées écrit les champs et crédite
        les contributeurs (celui qui est déjà crédité dans ce rôle est
        écarté). Les champs écrits comptent comme \emph{modifiés à la
        main} : un téléchargement automatique ultérieur des
        métadonnées ne les écrase jamais ;
  \item une proposition de couverture installe l'image comme
        couverture du titre, elle aussi protégée des téléchargements
        automatiques de couverture.
\end{itemize}

Le message de succès renvoie à la série ou au titre.

\subsection{Propositions périmées}

Une proposition est vérifiée à nouveau au moment où vous l'acceptez.
Si son objet a changé depuis le dépôt --- le titre ou la série a été
modifié ou supprimé, une position a été prise entre-temps, une série
du même nom existe désormais, la couverture a changé, un rôle de
contributeur a été supprimé ---,
\textbf{rien n'est appliqué} et la proposition devient
\emph{Périmée} ; le message explique pourquoi. Demandez à l'agent une
nouvelle proposition fondée sur le catalogue tel qu'il est
aujourd'hui, ou faites la modification à la main.

Seul compte ce que la proposition touche. Accepter une couverture ne
rend pas périmée une proposition de métadonnées en attente pour le
même titre, et accepter une description ne rend pas périmée une
proposition d'éditeur : décidez-les dans l'ordre que vous voulez. Une
proposition de métadonnées devient périmée quand l'un de \emph{ses}
champs a changé depuis son dépôt (à la main, par une autre proposition
acceptée ou par un téléchargement de métadonnées), une proposition de
couverture quand la couverture du titre a changé --- et les deux quand
le titre lui-même a été supprimé ou ré-identifié à la main (son nom,
son ISBN ou son type de média corrigé ; une correction de titre acceptée
depuis une proposition laisse les autres propositions valables) ou qu'un rôle de contributeur proposé a été
supprimé. Modifier un champ que la proposition ne touche pas ne rend
rien périmé.

Révoquer une clé (section~\ref{sec:agents-revoke}) rend aussi
périmées toutes ses propositions en attente ; la page de détail
indique alors qu'elles ont été décidées par \emph{le système}.

\subsection{Refuser}

\textbf{Refuser} clôt la proposition sans rien changer. Vous pouvez
donner un \textbf{motif} (facultatif, 1000 caractères au plus) :
l'agent le lit dans son historique (\texttt{list\_my\_proposals}) et
peut en tirer parti. Une couverture proposée est supprimée quand sa
proposition est refusée.

\subsection{Le délai avant nouvelle soumission}
\label{sec:agents-delay}

Après un refus, une proposition \emph{équivalente} est refusée pendant
un certain temps, pour qu'un agent ne dépose pas la même chose encore
et encore. Équivalente signifie de même type et de même objet : la
même série (une nouvelle série par son nom, sans tenir compte de la
casse ni des accents) ; pour des métadonnées, le même titre et
\emph{l'un quelconque} des champs refusés --- un éditeur refusé bloque
toute proposition qui comprend l'éditeur, l'année restant libre ; pour
une couverture, le même titre \emph{et la même image} --- une autre
image peut être proposée tout de suite. La
règle vaut pour toutes les clés : créer une nouvelle clé ne la remet
pas à zéro.

Le délai se règle dans \textbf{Administration \(\rightarrow\) Système
\(\rightarrow\) Propositions}\index{délai avant nouvelle soumission}
(section~\ref{sec:system}) : de 0 (aucun délai) à 3650 jours,
\textbf{90 jours} par défaut. Un changement s'applique à la
proposition suivante, sans redémarrage.

\section{Tomes attendus et wish list}
\label{sec:expected-tomes}

Un \emph{tome attendu}\index{tome attendu} nomme un tome d'une série
que vous ne possédez pas encore : une position, un titre et, s'il est
connu, un ISBN. Les tomes attendus viennent de propositions de série
acceptées, ou vous les saisissez vous-même.

\begin{itemize}
  \item \textbf{Sur la page de la série}, un tome attendu occupe la
        case de sa lacune avec son titre, son ISBN et la mention
        \emph{Attendu}. Tout le monde le voit ; les compteurs de
        tomes possédés et manquants ne changent pas, puisque le tome
        manque toujours.
  \item \textbf{En saisir un à la main} (Bibliothécaire ou Admin) :
        \textbf{Ajouter un tome attendu} dans les actions de la
        série, ou \textbf{Nommer ce tome} sur une lacune vide.
        \textbf{Modifier} et \textbf{Supprimer} apparaissent sous
        chaque tome.
  \item \textbf{Ajouter à la wish list} inscrit le tome sur votre
        wish list (section~\ref{sec:wishlist}) avec son titre, son
        ISBN et une note qui nomme la série et la position. Un double
        clic ne crée pas de seconde entrée.
  \item \textbf{Quand vous cataloguez le livre} --- en scannant son
        ISBN, ou en créant le titre à la main avec cet ISBN ---,
        mybibli reconnaît le tome attendu, place le nouveau titre à
        cette position de la série, et la lacune disparaît. Si le
        livre figurait sur votre wish list, le bandeau habituel
        propose de l'en retirer (section~\ref{sec:wishlist}). Si la
        position est entre-temps occupée par un autre titre, rien
        n'est rattaché et un avertissement indique quoi vérifier.
  \item \textbf{Tomes attendus à revoir} --- un tome qui ne tient plus
        dans une lacune (le nombre de volumes de la série a été
        abaissé sous sa position, ou un titre occupe désormais sa
        position) est listé sous la grille avec la raison. Changez sa
        position ou supprimez-le.
\end{itemize}

\section{Couvertures proposées : la zone d'attente}
\label{sec:agents-staging}

Une couverture proposée est contrôlée comme une couverture que vous
téléversez vous-même, puis réduite à 400\,px de large au plus
(JPEG) --- ce que vous examinez est exactement ce que l'acceptation
installe. Jusqu'à votre décision, elle attend dans un dossier caché
\texttt{.staging}\index{couvertures!.staging} \emph{à l'intérieur du
dossier des couvertures} (\texttt{COVERS\_DIR} ; sous Docker,
\texttt{/app/covers/.staging} sur le volume
\texttt{mybibli\_covers}). Conséquences :

\begin{itemize}
  \item Rien à configurer : toute installation qui conserve ses
        couvertures conserve aussi les couvertures proposées d'une
        mise à jour à l'autre.
  \item Une image proposée n'est jamais publique : l'adresse
        \texttt{/covers} refuse les dossiers cachés, et la page
        Propositions la sert aux seuls bibliothécaires, tant que la
        proposition est en attente.
  \item Le nettoyage est automatique. Refuser une proposition supprime
        son image ; l'accepter la met en place ; révoquer une clé
        supprime les images de ses propositions en attente. La purge
        planifiée (section~\ref{sec:system}, \emph{intervalle d'auto-purge}) retire tout reste à son passage
        suivant --- l'image d'une proposition devenue périmée, par
        exemple ; un dépôt jamais achevé seulement après une heure.
  \item Les sauvegardes du volume des couvertures
        (chapitre~\ref{ch:backup}) incluent ce dossier ; le restaurer
        est sans danger, la purge fait le tri.
\end{itemize}

\section{Révoquer la clé d'un agent}
\label{sec:agents-revoke}

Dans \textbf{Administration \(\rightarrow\) Clés d'API}, révoquez la
clé. L'agent est refusé dès sa requête suivante, toutes ses
propositions en attente deviennent périmées et leurs couvertures
proposées sont supprimées. Ses propositions acceptées et refusées
restent dans l'historique. Pour réadmettre l'agent, créez une nouvelle
clé et mettez à jour la configuration du client.
